\documentclass{beamer}

\input{preamble.tex}

\title{Instrumental Variables}
\subtitle{Session 1: Foundations and Mechanics}
\author{Jingxuan Du \\ \smallskip {\small Dongbei University of Finance and Economics (DUFE)} \\ \smallskip {\small \href{mailto:dujingxuan6@dufe.edu.cn}{dujingxuan6@dufe.edu.cn}}}
\date{Empirical Methods for PhD Students}

\begin{document}

\imageframe{./lecture_includes/cover_intro.png}

\section{Introductions}

\begin{frame}{Who Am I?}

A Professor of Economics at Dongbei University of Finance and Economics (DUFE)\pause

A researcher working on empirical methods in economics\pause
\begin{itemize}
  \item Applied econometrics with a focus on causal identification
  \item IV-based analyses in labor, trade, and development economics
  \item Policy evaluation using quasi-experimental methods
\end{itemize}\pause

Committed to teaching rigorous, practical empirical tools to the next generation of economists

\end{frame}

\begin{frame}{Why IV? Motivation}

Economists care about \emph{causal} effects, not just correlations\pause
\begin{itemize}
  \item Does education raise earnings, or do smarter people get more education?
  \item Does trade competition hurt employment, or do weak industries seek protection?
  \item Does infrastructure investment drive growth, or do wealthy regions get more roads?
\end{itemize}\medskip\pause

The fundamental problem: \textbf{selection bias}
\begin{itemize}
  \item Treatment assignment is rarely random in observational data
  \item OLS captures both causal effects and selection differences
\end{itemize}\medskip\pause

\textbf{Instrumental variables (IV)} offer a principled solution: find a source of exogenous variation in the treatment that is unrelated to the unobserved confounders

\end{frame}

\begin{frame}{What Is This PhD Course?}

A two-session intensive on IV methods, aimed at PhD students in economics\pause

\begin{itemize}
  \item Session 1 (today): IV mechanics, identification, and foundational applications
  \item Session 2: IV interpretation (LATE), formula instruments, and advanced topics
\end{itemize}\pause\medskip

Emphasis on \emph{practice}: IV is a tool meant to be used in your own research!\pause

\begin{itemize}
  \item Please ask questions throughout
  \item We will connect methods to published empirical papers
\end{itemize}\pause\medskip

Featured applications: seminal papers on \textbf{Chinese economic topics} from top journals (AER, JPE, ReStud)

\end{frame}

\begin{frame}{Session 1 Roadmap}

\begin{enumerate}
  \item \textbf{IV Mechanics}
  \begin{itemize}
    \item Just-identified IV: identification, estimation, inference
    \item Overidentification and 2SLS
    \item Weak instruments and many-IV bias
  \end{itemize}\medskip
  \item \textbf{Chinese Economics Applications: Session I}
  \begin{itemize}
    \item Paper 1: The China Syndrome (Autor, Dorn, Hanson 2013, AER)
    \item Paper 2: Roads and Growth in China (Banerjee, Duflo, Qian 2020, JPE)
  \end{itemize}\medskip
  \item \textbf{Conclusions and Preview of Session 2}
\end{enumerate}

\end{frame}

\section{IV Mechanics}

\subsection{Just-Identified IV}

\begin{frame}{Three Requirements for a Valid IV}

A valid instrument $Z_i$ must satisfy three conditions:\medskip\pause

\begin{enumerate}
  \item \textbf{Relevance (First Stage):} $Z_i$ must be correlated with the endogenous treatment $D_i$
  \begin{itemize}
    \item Testable: check the first-stage F-statistic (rule of thumb: $F > 10$)
  \end{itemize}\medskip\pause
  \item \textbf{As-Good-As-Random Assignment (Independence):} $Z_i$ must be uncorrelated with unobserved confounders $\varepsilon_i$
  \begin{itemize}
    \item Not directly testable; requires economic reasoning and robustness checks
  \end{itemize}\medskip\pause
  \item \textbf{Exclusion Restriction:} $Z_i$ affects outcome $Y_i$ \emph{only} through its effect on $D_i$
  \begin{itemize}
    \item No direct effect of $Z_i$ on $Y_i$; testable indirectly via falsification
  \end{itemize}
\end{enumerate}\medskip\pause

In Session 2, we add a 4th: \textbf{Monotonicity} (for LATE interpretation)

\end{frame}

\begin{frame}{Preliminaries: Parameters, Estimands, and Estimators}
\vspace{-0.2cm}
Three distinct objects, not always clearly distinguished: \pause
{\small
\begin{itemize}
  \item \bgVioletRed{Parameters} come from economic (or other) models of the world
  \begin{itemize}\itemsep0em
    \item E.g. a ``structural'' model of supply and demand, or a potential outcome model relating schooling to earnings
    \item They set the target for empirical analyses: what we want to know
  \end{itemize}\pause
  
  \item \bgPictonBlue{Estimands} are functions of the population data distribution 
  \begin{itemize}\itemsep0em
    \item E.g. a difference in means or ratio of population regression coef's
    \item We make assumptions to link parameters \& estimands (identification)
  \end{itemize}\pause
  
  \item \bgSun{Estimators} are functions of observed data (i.e. the ``sample'')
  \begin{itemize}\itemsep0em
    \item E.g. a difference in sample means or ratio of OLS coefficients 
    \item Since data are random, so are estimators. Each has a distribution
    \item We use knowledge of estimator distributions to learn about estimands (inference) and thus identified parameters
  \end{itemize}
\end{itemize}
}
\end{frame}

\begin{frame}{The Lay of the Land}
  \begin{adjustbox}{width = 0.9\textwidth, center}
    \begin{tikzpicture}
      \node[violet-red] (theory) at (-3.5, 1.25) {Economic Theory};
      \node[state,rectangle,violet-red] (parameters) at (-3.5, 0) {Parameters};
      
      \node[picton-blue] (distribution) at (0, 1.25) {Data Distribution};
      \node[state,rectangle,picton-blue] (estimands) at (0, 0) {Estimands};

      \node[sun] (data) at (3.5, 1.25) {Observed Datasets};
      \node[state,rectangle,sun] (estimators) at (3.5, 0) {Estimators};

      \path[violet-red, thick, shorten >=4pt] (theory) edge (parameters);
      \path[picton-blue, thick, shorten >=4pt] (distribution) edge (estimands);
      \path[sun, thick, shorten >=4pt] (data) edge (estimators);

      \path[thick, shorten >=4pt, shorten <=4pt] (estimands) edge[bend left= 40] node[below, el, yshift=-4pt] {Identification} (parameters);
      \path[thick, shorten >=4pt, shorten <=4pt] (estimators) edge[bend left= 40] node[below, el, yshift=-4pt] {Statistical Inference} (estimands);
    \end{tikzpicture}
  \end{adjustbox}

  \vspace{10mm}
  Separating out the very different tasks of identification and inference can help make our lives easier\smallskip\pause{}
\begin{itemize}
\item Today we'll start with the mechanics of IV estimands, talk briefly about estimation, then spend some time talking about identification
\end{itemize}
\end{frame}

\begin{frame}{An Example}
Human capital theory tells us that taking econometrics courses will likely boost later-life productivity\pause
\begin{itemize}
  \item Parameter: returns to enrolling in this course $\violetRed{\beta}$, measured in some outcome $Y_i$ (e.g. lifetime earnings or \# of top publications)
  \item Simple causal model: $Y_i= \violetRed{\beta} D_i + \varepsilon_i$, where $D_i \in \{0,1\}$ indicates enrollment and $\varepsilon_i$ is $i$'s outcome without the course
\end{itemize}\pause\medskip

We observe $Y_i$ and $D_i$ for some sample of individuals $i=1,\dots,N$
\begin{itemize}
  \item We fire up Stata and \code{reg Y D, r}. \pause{} How do we interpret the results?
\end{itemize}

\end{frame}

\begin{frame}{Population Regression and Endogeneity}
In large samples ($N\rightarrow\infty$), the OLS \emph{estimator} $\sun{\widehat{\beta}^{OLS}}$ gets arbitrarily close to (consistently estimates) the regression \emph{estimand} $\pictonBlue{\beta^{OLS}}$\pause
\begin{itemize} 
\item The \emph{identification} question is thus whether or not $\pictonBlue{\beta^{OLS}}=\violetRed{\beta}$
\end{itemize}\medskip\pause{}
By definition, $\pictonBlue{\beta^{OLS}}=\frac{Cov(Y_i,D_i)}{Var(D_i)}$.\pause{} Plugging in our causal model:
\begin{align*}
\pictonBlue{\beta^{OLS}} &= \frac{Cov( \violetRed{\beta} D_i + \varepsilon_i,D_i)}{Var(D_i)}\\ \pause{}
&=\violetRed{\beta} + \frac{Cov(\varepsilon_i,D_i)}{Var(D_i)}
\end{align*}\pause{}
Thus, we have (regression) identification if and only if $Cov(\varepsilon_i,D_i)=0$\smallskip
\begin{itemize}\pause{}
\item Selection bias: people with certain potential outcomes $\varepsilon_i$ are more/less likely to take this course, such that $Cov(\varepsilon_i,D_i)\neq 0$
\end{itemize}
\end{frame}

\begin{frame}{Regression ``Exogeneity''}
  \begin{adjustbox}{width = 0.6\textwidth, center}
    \begin{tikzpicture}
      \node[state] (D) at (0,0) {$D$};
      \node[state] (Y) at (2,0) {$Y$};
      \node[state, ruby] (epsilon) at (1,1.5) {$\epsilon$};
    
      \path[violet-red, thick] (D) edge node[above, el] {$\violetRed{\beta}$} (Y);
      \path[ruby, thick] (epsilon) edge (Y);
    \end{tikzpicture}
  \end{adjustbox}
\end{frame}

\begin{frame}{Regression ``Endogeneity''}
  \begin{adjustbox}{width = 0.6\textwidth, center}
    \begin{tikzpicture}
      \node[state] (D) at (0,0) {$D$};
      \node[state] (Y) at (2,0) {$Y$};
      \node[state, ruby] (epsilon) at (1,1.5) {$\epsilon$};
    
      \path[violet-red, thick] (D) edge node[above, el] {$\violetRed{\beta}$} (Y);
      \path[ruby, thick] (epsilon) edge (Y);
      \path[ruby, thick] (epsilon) edge (D);
    \end{tikzpicture}
  \end{adjustbox}
\end{frame}

\begin{frame}{The IV Solution}

Suppose this course was ``oversubscribed'' (more people wanted to attend than space allowed), so seats were determined by lottery
\begin{itemize}
\item $Z_i\in\{0,1\}$ indicates randomized admission to the course\pause{}
\item Randomness + no direct effects of $Z_i$ on $Y_i$ implies $Cov(Z_i,\varepsilon_i)=0$
\end{itemize}\medskip\pause{}

Plugging in the model for $\varepsilon_i=Y_i-\violetRed{\beta} D_i$, we now have IV identification:
\begin{align*}
Cov(Z_i,Y_i-\violetRed{\beta} D_i)=0 \pause{} \implies \violetRed{\beta} = \frac{Cov(Z_i,Y_i)}{Cov(Z_i,D_i)}\pause{}\equiv \pictonBlue{\beta^{IV}},
\end{align*}
so long as $Cov(Z_i,D_i)\neq 0$.\pause{} We can estimate this by $\sun{\widehat{\beta}^{IV}}=\frac{\widehat{Cov}(Z_i,Y_i)}{\widehat{Cov}(Z_i,D_i)}$\pause{}
\begin{itemize}
\item Or, in Stata, \code{ivreg2 Y (D=Z), r}
\end{itemize}

\end{frame}

\begin{frame}{The IV Solution}
  \begin{adjustbox}{width = 0.7\textwidth, center}
    \begin{tikzpicture}
      \node[state] (D) at (0,0) {$D$};
      \node[state] (Y) at (2,0) {$Y$};
      \node[state, ruby] (epsilon) at (1,1.5) {$\epsilon$};
      \node[state, kelly] (Z) at (-2,0) {$Z$};
    
      \path[violet-red, thick] (D) edge node[above, el] {$\violetRed{\beta}$} (Y);
      \path[ruby, thick] (epsilon) edge (Y);
      \path[ruby, thick] (epsilon) edge (D);
      \path[kelly, thick] (Z) edge (D);
    \end{tikzpicture}
  \end{adjustbox}
\bigskip

Note: no arrow connecting $\varepsilon$ and $Z$ (``as-good-as-random assignment''), and no arrow from $Z$ to $Y$ directly (``exclusion''). We'll come back to both
\end{frame}

\begin{frame}{Reduced Form and First Stage}

We're usually pretty comfortable w/regression; IV feels more mysterious\smallskip
\begin{itemize}
\item But ``one weird trick'' links the IV estimand back to regression:
\begin{align*}
\pictonBlue{\beta^{IV}} = \frac{Cov(Z_i,Y_i)}{Cov(Z_i,D_i)} \pause = \frac{Cov(Z_i,Y_i)/Var(Z_i)}{Cov(Z_i,D_i)/Var(Z_i)} \pause \equiv \frac{\pictonBlue{\rho}}{\pictonBlue{\pi}}
\end{align*}
where $\pictonBlue{\rho}$ and $\pictonBlue{\pi}$ come from two population regressions:
\begin{align*}
Y_i &= \kappa+ \pictonBlue{\rho} Z_i + \nu_i \hspace{0.2cm}\text{The ``reduced form''}\\
D_i &= \mu + \pictonBlue{\pi} Z_i + \eta_i\hspace{0.2cm}\text{The ``first stage''}
\end{align*}
\end{itemize}

\end{frame}

\begin{frame}{Angrist (1990): Draft Lottery IV}

Angrist famously used Vietnam-era draft eligibility as an instrument to estimate the earnings effects of military service 
\begin{itemize}
\item Let $Z_i\in\{0,1\}$ be an indicator for draft eligibility, $D_i\in\{0,1\}$ be an indicator for military service, and $Y_i$ measure later-life earnings
\end{itemize}\medskip\pause

Here $\pictonBlue{\beta^{IV}} = \frac{Cov(Z_i,Y_i)/Var(Z_i)}{Cov(Z_i,D_i)/Var(Z_i)} = \frac{E[Y_i\mid Z_i=1]-E[Y_i\mid Z_i=0]}{E[D_i\mid Z_i=1]-E[D_i\mid Z_i=0]}$ has a special name, because $Z_i$ is binary: the \bgPictonBlue{Wald estimand}\pause{}
\begin{itemize}
\item First stage $E[D_i\mid Z_i=1]-E[D_i\mid Z_i=0]$: effect of eligibility on the \emph{probability} of military service (b/c $D_i$ is binary)
\item Reduced form $E[Y_i\mid Z_i=1]-E[Y_i\mid Z_i=0]$: effect of eligibility on adult earnings (measured in 1971, 1981...)
\end{itemize}\medskip\pause

IV interprets the latter causal effect in terms of the former
\end{frame}

\imageframe{./lecture_includes/angrist_1990.png}

\begin{frame}{Adding Controls}

We might only think our $Z_i$ is exogenous controlling for some vector $W_i$
\begin{itemize}
\item Just add controls to the reduced form and first stage! $\beta^{IV}=\frac{\rho}{\pi}$ for
\vspace{-0.3cm}
\begin{align*}
Y_i &= \kappa+ \rho Z_i + W_i^\prime\phi +  \nu_i \\
D_i &= \mu +\pi Z_i + W_i^\prime\psi + \eta_i
\end{align*}

\vspace{-0.2cm}
with the estimator $\hat\beta^{IV}$ defined analogously \pause{}\medskip
\item The Frisch-Waugh-Lovell theorem tells us the effective instrument is now $\tilde{Z_i}$: the residuals from regressing $Z_i$ on $W_i$ in the population\pause{}\medskip
\item E.g. if $W_i$ is a vector of dummies for randomization strata in an RCT, then $\tilde{Z_i}$ captures the within-strata variation in $Z_i$
\end{itemize}
\end{frame}

\begin{frame}{Multiple Treatments}

We might be interested in a multi-dimensional model: $Y_i=X_i^\prime\beta+\varepsilon_i$
\begin{itemize}
\item Instrument vector $Z_i$, with $dim(Z_i)=dim(X_i)$ (``just-identified'')
\item Population residuals $\tilde{Z}_i$, given some vector of controls $W_i$
\end{itemize}\medskip\pause{}
Suppose $Cov(\tilde{Z}_i,\varepsilon_i)=0$. Then, just as before, we have identification:
\begin{align*}
Cov(\tilde{Z}_i,Y_i-X_i^\prime\beta)=0 \implies \beta = Cov(\tilde{Z}_i,X_i)^{-1}Cov(\tilde{Z}_i,Y_i)\equiv \beta^{IV},
\end{align*}
so long as $Cov(\tilde{Z}_i,X_i)$ is full-rank.\pause{} Equivalently, $\beta^{IV}=\pi^{-1}\rho$ where:
\vspace{-0.2cm}
\begin{align*}
Y_i&=Z_i^\prime\rho + W_i^\prime\phi+\nu_i\\
X_i&=\pi Z_i+W_i^\prime\psi_i+\eta_i,
\end{align*}

\vspace{-0.3cm}
with estimation working just as you'd think
\end{frame}

\subsection{Overidentification}

\begin{frame}{Multiple Instruments}

What happens when $dim(Z_i)=L>J=dim(X_i)$? \emph{Overidentification}:
\vspace{-0.2cm}
\begin{align*}
Cov(\tilde{Z}_i,Y_i-X_i^\prime\beta)=0 \implies \underbrace{Cov(\tilde{Z}_i,Y_i)}_{L\times 1}=\underbrace{Cov(\tilde{Z}_i,X_i^\prime)}_{L\times J}\underbrace{\beta}_{J\times 1}
\end{align*}
\vspace{-0.5cm}
\begin{itemize}
\item We can drop any $L-J$ instruments to identify $\beta$\pause{}
\item More generally, we can take any full-rank linear combination $\tilde{Z}_i^*=M\tilde{Z}_i$ for $J\times L$ matrix $M$ such that $Cov(\tilde{Z}_i^*,X_i)$ is invertible
\end{itemize}
\vspace{-0.3cm}
\begin{align*}
&\underbrace{M\cdot Cov (\tilde{Z}_i,Y_i)}_{Cov(\tilde{Z}_i^*,Y_i)}=\underbrace{M\cdot Cov(\tilde{Z}_i,X_i^\prime)}_{Cov(\tilde{Z}_i^*,X_i^\prime)}\beta\\
&\implies \beta = \left(M\cdot Cov(\tilde{Z}_i,X_i^\prime)\right)^{-1} M\cdot Cov(\tilde{Z}_i,Y_i)
\end{align*}\pause{}
This defines a \emph{class} of IV estimands/estimators, indexed by $M$
\end{frame}

\begin{frame}{Two-Stage Least Squares (2SLS)}

2SLS sets $M^\prime=\pi$: the matrix of first-stage coefficients
\begin{itemize}
\item This makes $\tilde{Z}_i^*=\tilde{Z}_i^\prime\pi$ the (residualized) first-stage fitted values 
\item Intuitively: combine IVs according to their predictiveness of $X_i$; \\ we should expect this to decrease the IV standard error
\end{itemize}\bigskip\pause{}
Since the first-stage from regressing $X_i$ on this $\tilde{Z}_i^*$ is one (by construction), $\beta^{2SLS}$ can be obtained in two \emph{stages}:
\begin{enumerate}
\item Regress $X_i$ on $Z_i$ and $W_i$ (first stage)
\item Regress $Y_i$ on first-stage fitted values and $W_i$ (second stage)
\end{enumerate}\bigskip\pause{}
Note: here we're talking about the 2SLS \emph{estimand}, but the exact same logic holds for the 2SLS \emph{estimator} (i.e. two stages of OLS)
\begin{itemize}
\item That said, you should never run 2SLS by hand. Let \emph{ivreg2} do it!
\end{itemize}
\end{frame}

\begin{frame}{Angrist-Krueger '91: The Power of Overidentification}

\includegraphics[scale=0.5]{./lecture_includes/ak_table.png}

\end{frame}

\begin{frame}{2SLS is a Many-Splendored Thing}

There is another (I think more useful) way to understand 2SLS: \\ as a weighted average of just-identified IVs:
\begin{align*}
\beta^{2SLS}&=\left(\pi^\prime Cov(\tilde{Z}_i,X_i^\prime)\right)^{-1} \pi^\prime Cov(\tilde{Z}_i,Y_i)\\
&=(\pi^\prime Var(\tilde{Z}_i)\pi)^{-1}\pi^\prime Var(\tilde{Z}_i)\rho,
\end{align*}
This is the formula for a $Var(\tilde{Z}_i)$-weighted regression of reduced-form coefficients $\rho$ on first-stage coefficients $\pi$ (with no constant)\pause{}
\begin{itemize}
\item When $J=1$ (one treatment), this becomes $\beta^{2SLS}=\sum_\ell \omega_\ell \beta_\ell^{IV}$ where $\omega_\ell=(\pi^\prime Var(\tilde{Z}_i)\pi)^{-1}\pi^\prime Var(\tilde{Z}_i)^\prime_\ell\pi_\ell$ and $\beta^{IV}_\ell = \rho_\ell/\pi_\ell$
\item Intuitively: 2SLS combines multiple ``one-at-a-time'' IVs $\beta_\ell^{IV}$
\end{itemize}

\end{frame}

\begin{frame}{Angrist-Hull '23: ``Visual IV'' for Cancer Screening Trials}
\begin{center}
\includegraphics[scale=0.5]{./lecture_includes/pragmaticTrials.png}
\vspace{-0.3cm}
\end{center}
Each dot gives a ($\rho_\ell,\pi_\ell$) for a trial $\ell$ where  randomized screening offers $Z_i$ instrument for screening participation $D_i$
\begin{itemize}
\item Slope of the weighted line-of-best fit through zero = 2SLS estimate
\end{itemize}
\end{frame}

\begin{frame}{Overidentification Tests}

Under the constant-effects causal model of $Y_i=X_i^\prime\beta+\varepsilon_i$, overidentification gives a way to test instrument validity
\begin{itemize}
\item All just-identified IVs should coincide: i.e. $\beta_\ell^{IV}=\beta$ for all $\ell$
\item Graphically: the $R^2$ from visual IV plots should $= 1$ in large samples
\end{itemize}\medskip\pause{}
Stata will automatically give you the $p$-value for this test when $L>J$
\begin{itemize}
\item If $p>0.05$, it means your $\hat{\beta}_\ell^{IV}$'s are all pretty similar to each other
\end{itemize}\medskip\pause{}
Don't place too much stock in overidentification tests, however:
\begin{itemize}
\item They tend to have low power (b/c individual $\hat{\beta}_\ell^{IV}$ tend to be noisy)
\item If they reject, it need not mean your instruments are invalid \\ (b/c of treatment effect heterogeneity -- more on this in Session 2)
\item Rejection doesn't tell you which IV is invalid (they all might be!)
\end{itemize}
\end{frame}

\subsection{Weak vs.\ Many-Weak Bias}

\begin{frame}{Weak Instruments}
\vspace{-0.2cm}

When running just-identified IV, people worry about instrument ``strength'' \vspace{-0.7cm}
\begin{itemize}
  \item Specifically the first stage F-statistic, which tests if $\pi=0$ 
\end{itemize}
\smallskip\pause

If $\pi$ is small relative to its standard error, we say the IV is ``weak''
\begin{itemize}
  \item Typically use the rule-of-thumb of $F<10$ (Staiger and Stock 1997)
  \item In this case the second-stage SEs will be large and the 2SLS estimate will tend to be biased towards the corresponding OLS
\end{itemize}\pause

Much has been made of this over the years, but Angrist and Koles\'{a}r (2022) show that we shouldn't worry too much
\begin{itemize}
  \item The SE increase tends to be large enough to ``cover up'' the bias, \\ so you're unlikely to reject the null of $\beta=0$ spuriously 
\end{itemize}

\end{frame}

\begin{frame}{Weak Instruments: Visualized}
\vspace{-0.2cm}
Monte Carlo: $Y_i=0\cdot D_i + \varepsilon_i$, $D_i=\pi Z_i+\eta_i$: $\pi=Var(\varepsilon_i)=Var(\eta_i)=1$
\begin{center}
\includegraphics[scale=0.45]{./lecture_includes/strongpi.png}
\end{center}

\end{frame}

\begin{frame}{Weak Instruments: Visualized}
\vspace{-0.2cm}
Monte Carlo: $Y_i=0\cdot D_i + \varepsilon_i$, $D_i=\pi Z_i+\eta_i$: $\pi=0.1$ (Weaker)
\begin{center}
\includegraphics[scale=0.35]{./lecture_includes/medpi.png}
\end{center}

\end{frame}

\begin{frame}{Weak Instruments: Visualized}
\vspace{-0.2cm}
Monte Carlo: $Y_i=0\cdot D_i + \varepsilon_i$, $D_i=\pi Z_i+\eta_i$: $\pi=0.01$ (Very Weak)
\begin{center}
\includegraphics[scale=0.35]{./lecture_includes/weakpi.png}
\end{center}

\end{frame}

\begin{frame}{Many IVs}
\vspace{-0.2cm}
A thornier problem is many-weak bias, when overidentified\smallskip
\begin{itemize}
\item This also tends to manifest in low first-stage F's, and also causes 2SLS to be biased towards OLS
\end{itemize}\bigskip\pause{}

Unlike when just-id., however, with many weak IVs the SE's go \emph{down}\smallskip
\begin{itemize}
  \item Intuitively, a more flexible FS tends to fit $D_i$ better $\rightarrow$ more power\smallskip
  \item But we can have \emph{overfitting} with lots of instruments, which essentially recreates the (endogenous) variation in $D_i$ 
\end{itemize}\bigskip\pause{}

This became a high-profile problem with Angrist-Krueger '91, where the QOB instrument was interacted with many state/year FEs
\smallskip
\begin{itemize}
  \item These days folks don't make this mistake ... but many-IV bias can be
lurking in other settings with constructed instruments (e.g. judge IV)
\end{itemize}

\end{frame}

\begin{frame}{Many Instruments: Visualized}
\vspace{-0.2cm}
$Y_i=0\cdot D_i + \varepsilon_i$, $D_i=\pi Z_{i1}+\sum_{\ell>1} 0\cdot Z_{i\ell}+\eta_i$: IV w/ $Z_{i1}$ only
\begin{center}
\includegraphics[scale=0.35]{./lecture_includes/fewz.png}
\end{center}

\end{frame}

\begin{frame}{Many Instruments: Visualized}
\vspace{-0.2cm}
$Y_i=0\cdot D_i + \varepsilon_i$, $D_i=\pi Z_{i1}+\sum_{\ell>1} 0\cdot Z_{i\ell}+\eta_i$: IV w/ $Z_{i1},\dots,Z_{i10}$
\begin{center}
\includegraphics[scale=0.35]{./lecture_includes/somez.png}
\end{center}

\end{frame}

\begin{frame}{Many Instruments: Visualized}
\vspace{-0.2cm}
$Y_i=0\cdot D_i + \varepsilon_i$, $D_i=\pi Z_{i1}+\sum_{\ell>1} 0\cdot Z_{i\ell}+\eta_i$: IV w/ $Z_{i1},\dots,Z_{i100}$
\begin{center}
\includegraphics[scale=0.35]{./lecture_includes/manyz.png}
\end{center}

\end{frame}

\begin{frame}{What to Do?}
\vspace{-0.2cm}
Aim for few instruments, and check your F's after every \emph{ivreg}\smallskip
\begin{itemize}
\item State of the art: Montiel Olea and Pflueger '15; \code{weakivtest} in Stata\smallskip
\item Staiger-Stock rule-of-thumb ($F>10$) still seems widely held\smallskip
\item See Lee et al. (2020) and Keane and Neal (2022) for some discussions of additional subtleties
\end{itemize}
\medskip\pause{}

If your F is small, some things to consider:\smallskip
\begin{itemize}
\item Is there a better functional form for your instrument?\smallskip
\item Do interactions with covariates help? (note: beware many-weak!)\smallskip
\item Does changing the covariate set help? (note: beware invalidity!)\smallskip
\item Check results w/a more robust approach (e.g. Anderson-Rubin, JIVE)
\end{itemize}

\end{frame}

\section{Chinese Economic Applications: Session I}

\begin{frame}{IV Methods in Chinese Economic Research}

China provides rich empirical settings for IV research:\pause
\begin{itemize}
  \item Rapid trade liberalization $\rightarrow$ natural experiments in import competition
  \item Massive infrastructure investment $\rightarrow$ exogenous variation in connectivity
  \item Historical institutions and policies $\rightarrow$ long-run IV strategies
  \item Large administrative datasets at province, county, and firm level
\end{itemize}\medskip\pause

In this section, we examine two landmark papers using IV to study Chinese economic phenomena published in \textbf{top economics journals}:

\begin{enumerate}
  \item \textbf{The China Syndrome} (Autor, Dorn, Hanson 2013, \emph{AER}) \\
  {\small Shift-share IV to study effects of Chinese import competition on US labor markets}
  \item \textbf{On the Road} (Banerjee, Duflo, Qian 2020, \emph{JPE}) \\
  {\small Historical railroad plan as IV for infrastructure effects on Chinese economic growth}
\end{enumerate}

\end{frame}

\begin{frame}{Paper 1: ``The China Syndrome''}

{\small
\textbf{Citation:} Autor, D., Dorn, D., and Hanson, G.H. (2013). ``The China Syndrome: Local Labor Market Effects of Import Competition in the United States.'' \emph{American Economic Review}, 103(6): 2121--2168.
}\medskip\pause

\textbf{Context:} China's WTO accession (2001) accelerated its integration into world trade. Chinese imports to the US grew from \$100B in 2000 to over \$400B by 2007.\medskip\pause

\textbf{Central Question:} Does rising import competition from China cause adverse labor market outcomes in US commuting zones?
\begin{itemize}
  \item Challenge: Chinese imports may respond to demand conditions in the US that independently affect local employment --- \emph{reverse causality and omitted variable bias}
\end{itemize}

\end{frame}

\begin{frame}{ADH 2013: IV Strategy (Shift-Share)}

\textbf{The Endogeneity Problem:}
\begin{itemize}
  \item Commuting zones with weaker manufacturing may both import more from China \emph{and} have declining employment for other reasons
\end{itemize}\medskip\pause

\textbf{IV Solution --- Shift-Share Instrument:}
\begin{align*}
Z_i = \sum_k \underbrace{s_{ik}}_{\text{share}} \times \underbrace{g_k^{OTH}}_{\text{shock}}
\end{align*}
\begin{itemize}
  \item \textbf{Shares} $s_{ik}$: initial (1990) employment share of industry $k$ in CZ $i$
  \item \textbf{Shocks} $g_k^{OTH}$: growth of Chinese imports in \emph{8 other} high-income countries (Australia, Denmark, Finland, Germany, Japan, New Zealand, Spain, Switzerland)
\end{itemize}\medskip\pause

\textbf{Identification intuition:} Imports in other countries reflect Chinese productivity growth (the exogenous shock), but not US-specific demand shocks

\end{frame}

\begin{frame}{ADH 2013: Data and Setting}

\textbf{Unit of observation:} 722 US commuting zones (CZs), two periods (1990--2000, 2000--2007)\medskip\pause

\textbf{Outcome variables:}
\begin{itemize}
  \item Manufacturing employment share in working-age population
  \item Wages, hours worked, industry switching
  \item Unemployment insurance, disability insurance receipt
\end{itemize}\medskip\pause

\textbf{Data sources:}
\begin{itemize}
  \item US Decennial Census and American Community Survey (employment/wages)
  \item UN Comtrade (bilateral import flows by industry)
  \item County Business Patterns (establishment-level employment)
  \item Bureau of Labor Statistics (unemployment and transfer payments)
\end{itemize}

\end{frame}

\begin{frame}{ADH 2013: Original Results}

\begin{center}
\includegraphics[width=0.85\textwidth]{./lecture_includes/adh_bhj.png}
\end{center}

\end{frame}

\begin{frame}{ADH 2013: Refined Results (Borusyak, Hull, Jaravel 2022)}

\begin{center}
\includegraphics[width=0.85\textwidth]{./lecture_includes/adh_bhj_restud.png}
\end{center}
{\small
\begin{itemize}
\item Columns 3--7 include the key $\sum_k s_{ik}\times(\text{period FE})$ control
\item Standard errors / first-stage F-stats computed via \emph{ssaggregate}
\end{itemize}
}

\end{frame}

\begin{frame}{ADH 2013: Findings and Impact}

\textbf{Main findings:}
\begin{itemize}
  \item \$1,000 increase in Chinese import exposure per worker $\rightarrow$ $\approx 2$--$2.5\%$ decline in the manufacturing employment share in commuting zones\pause
  \item Estimated 2--2.4 million manufacturing jobs lost in 1990--2007 due to Chinese import competition\pause
  \item Significant effects on wages, disability insurance, and income transfers
  \item Effects persist over a decade --- labor markets adjust slowly
\end{itemize}\medskip\pause

\textbf{Why it matters for IV methods:}
\begin{itemize}
  \item Introduced and popularized the \textbf{shift-share IV} design in economics
  \item Borusyak, Hull, and Jaravel (2022) provide formal conditions for validity: key is exogeneity of \emph{shocks} ($g_k$), not shares
  \item Generated a large literature on SSIV refinements and diagnostics
\end{itemize}

\end{frame}

\begin{frame}{Paper 2: ``On the Road''}

{\small
\textbf{Citation:} Banerjee, A., Duflo, E., and Qian, N. (2020). ``On the Road: Access to Transportation Infrastructure and Economic Growth in China.'' \emph{Journal of Political Economy}, 128(8): 2967--3001.
}\medskip\pause

\textbf{Context:} China built one of the world's most extensive road and railroad networks during the 1990s--2000s. By 2010, China had the world's second-largest highway network.\medskip\pause

\textbf{Central Question:} Does access to transportation infrastructure (national highways and railroads) cause economic growth in Chinese counties?
\begin{itemize}
  \item Challenge: roads are built to serve economically active areas $\rightarrow$ strong \emph{reverse causality}: wealthier counties attract more infrastructure
  \item OLS estimates would be biased upward
\end{itemize}

\end{frame}

\begin{frame}{BDQ 2020: IV Strategy}

\textbf{Identification challenge:} Infrastructure placement is endogenous --- roads go where economic activity is or is expected to be.\medskip\pause

\textbf{IV Approach --- Historical Planned Network:}
\begin{itemize}
  \item Key idea: use the \textbf{straight-line distance} between county seats and the pre-existing (pre-reform era) trunk transportation network as an instrument
  \item The historical network layout was driven by \textbf{political and military} considerations, not economic development potential of specific localities
\end{itemize}\medskip\pause

\textbf{Instrument construction:}
\begin{enumerate}
  \item Identify major city pairs linked in the 1962 railroad plan
  \item Draw straight lines between them --- counties near these lines are more likely to receive highways/railroads later
  \item Closeness to historical straight-line routes $\perp$ current economic potential, conditional on baseline controls
\end{enumerate}

\end{frame}

\begin{frame}{BDQ 2020: Data and Methodology}

\textbf{Unit of observation:} 2,872 Chinese counties, 1986--2003\medskip\pause

\textbf{Outcome variables:}
\begin{itemize}
  \item Real GDP growth at the county level
  \item Industrial output, investment
  \item Nighttime light intensity (satellite proxy for economic activity)
\end{itemize}\medskip\pause

\textbf{Data sources:}
\begin{itemize}
  \item National Bureau of Statistics: county-level GDP and industrial output
  \item Ministry of Transportation: national highway and railroad GIS data
  \item Historical transport plan documents (1962 state-level railroad plan)
  \item County population census data for baseline controls
\end{itemize}\medskip\pause

\textbf{Specification:} Panel IV regression with county and year fixed effects, instrumenting actual access with proximity to historically planned lines

\end{frame}

\begin{frame}{BDQ 2020: Key Results and Lessons}

\textbf{Main findings:}
\begin{itemize}
  \item Proximity to a national highway increases GDP growth significantly
  \item Being within 10km of a highway $\rightarrow$ roughly 10--15\% higher annual GDP growth during the construction period\pause
  \item Effects are concentrated in manufacturing and trade-related sectors
  \item Counties with initially weaker connectivity benefit more from access
\end{itemize}\medskip\pause

\textbf{Why it matters for IV methods:}
\begin{itemize}
  \item Illustrates the ``\textbf{historical plan as instrument}'' strategy --- a common approach in development economics
  \item Distinguishes between actual placement (endogenous) and historically planned routes (exogenous to current economic conditions)
  \item Similar to Duranton and Turner (2011, AER) for US highways
\end{itemize}\medskip\pause

\textbf{Note:} Session 2 covers a related method --- \emph{Recentered IV} (Borusyak and Hull 2023) --- with China's High-Speed Rail as an illustration

\end{frame}

\begin{frame}{BDQ 2020: Heterogeneous Effects and Mechanisms}

\textbf{Heterogeneous effects:}\pause
\begin{itemize}
  \item Effects are \emph{larger} for counties that are initially more isolated (higher baseline travel costs)
  \item Effects are \emph{larger} for counties near the main infrastructure corridors (market access externalities)
\end{itemize}\medskip\pause

\textbf{Mechanisms:}
\begin{itemize}
  \item Reduced trade costs $\rightarrow$ easier access to input and output markets
  \item Increased firm entry and investment in newly connected counties
  \item Labor market integration: workers move toward opportunities
\end{itemize}\medskip\pause

\textbf{Validity checks:}
\begin{itemize}
  \item Placebo test: straight-line connections between \emph{non-major} cities do not predict economic outcomes
  \item Balance: counties near planned lines are not systematically different in pre-period trends
  \item First stage: strong first stage F-statistics across all specifications
\end{itemize}

\end{frame}

\begin{frame}{Comparing the Two Papers}

\begin{center}
\begin{tabular}{lll}
\hline
 & \textbf{ADH (2013)} & \textbf{BDQ (2020)} \\
\hline
Journal & AER & JPE \\
Setting & US CZs / China trade & Chinese counties \\
Treatment & Chinese import competition & Road/rail access \\
IV type & Shift-share & Historical plan \\
Exog. source & Imports in other countries & 1962 military plan \\
Unit & CZ $\times$ period & County $\times$ year \\
Key finding & Trade $\downarrow$ manufacturing & Infrastructure $\uparrow$ GDP \\
\hline
\end{tabular}
\end{center}\medskip\pause

\textbf{Common thread:} Both papers use a \textbf{``double-level''} IV logic: exogenous variation exists at a higher/different level (industries/planned routes), which is ``translated'' down to the unit of interest via a known formula

\end{frame}

\section{Conclusions}

\begin{frame}{Session 1 Summary}

\textbf{IV Mechanics:}
\begin{itemize}
  \item IV identification: $\beta^{IV} = \rho / \pi$ (reduced form over first stage)
  \item 2SLS with multiple instruments: efficiency-weighted combination
  \item Overidentification tests: valid under constant effects, interpret carefully
  \item Weak and many-IV bias: aim for $F > 10$ in first stage
\end{itemize}\medskip\pause

\textbf{Chinese Economics Applications:}
\begin{itemize}
  \item ADH (2013): Shift-share IV reveals large labor market costs of China trade shock
  \item BDQ (2020): Historical plan IV reveals positive infrastructure effects in China
  \item Both papers exemplify design-based causal inference with creative instruments
\end{itemize}

\end{frame}

\begin{frame}{Preview: Session 2}

In Session 2, we will cover:\medskip

\begin{enumerate}
  \item \textbf{IV Interpretation}
  \begin{itemize}
    \item What does IV identify? The Local Average Treatment Effect (LATE)
    \item Imbens-Angrist (1994): monotonicity and complier analysis
    \item Diff-in-diff and IV: connections and extensions
  \end{itemize}\medskip
  \item \textbf{Formula Instruments}
  \begin{itemize}
    \item Shift-share IV: formal conditions and implementation
    \item Recentered IV: beyond shift-share (Borusyak-Hull 2023)
    \item China's High-Speed Rail as a running illustration
  \end{itemize}\medskip
  \item \textbf{Chinese Applications: Session II}
  \begin{itemize}
    \item Meng, Qian, Yared (2015, ReStud): Causes of China's Great Famine
  \end{itemize}
\end{enumerate}

\end{frame}

\begin{frame}{Keep Calm and \code{ivreg2} On!}

\begin{center}
Thank you for your attention!

\bigskip
\textbf{Jingxuan Du} \\
Dongbei University of Finance and Economics \\
\smallskip
\url{dujingxuan6@dufe.edu.cn}

\bigskip
{\small Session 2 continues after a short break.}
\end{center}
\end{frame}

\end{document}
